%% kapa: exam=2 round=23 year=2012 period=1 subject=appraisal_practice subject_ko=감정평가실무 seq=1 number=1 points=40
%% source: Q-Net 감정평가사 기출문제 (cst00309) artlSeq=5056012 file=raw/1008/5056012/1교시(감정평가실무).hwp
%% conversion: kordoc 4.24.51 -> tools/gen2.py (auto); review=proofread-sonnet; flags=none
\begin{longquestion}{【문제 1】}
(주)A감정평가법인 甲감정평가사는 (주)K자산운용으로부터 감정평가 등을 의뢰 받았다. 주어진 자료를 활용하여 다음 물음에 답하시오. (40점) \par
\medskip\noindent 물음 1) 계약임대료를 기준으로 대상물건을 감정평가하시오. (10점) \par
\medskip\noindent 물음 2) 시장임대료를 기준으로 대상물건을 감정평가하시오. (15점) \par
\medskip\noindent 물음 3) 물음 1)의 감정평가액과 거래예정금액, 물음 2)의 감정평가액과 거래예정금액을 이용하여 각각의 순현재가치(NPV; Net Present Value)와 내부수익률(IRR; Internal Rate of Return)을 산출하시오. (5점) \par
\medskip\noindent 물음 4) 벤치마크 투자수익률을 내부수익률로 실현하기 위해 甲감정평가사가 (주)K자산운용에 제시할 대상물건의 거래예정금액을 결정하시오. (10점) \par
\textbf{<자료 1> 대상물건 개요} \par
1. 토지ㆍ건물 내역 \par
{\small\centering
\begin{tabular}{|c|c|c|}
\hline
\multirow{4}{*}{\makecell{토지}} & \makecell{소재지} & \makecell{서울특별시 G구 Y동 OO빌딩} \\
\cline{2-3}
 & \makecell{지 목} & \makecell{대} \\
\cline{2-3}
 & \makecell{면 적} & \makecell{2,833m\textsuperscript{2}} \\
\cline{2-3}
 & \makecell{용도지역} & \makecell{일반상업지역} \\
\hline
\multirow{8}{*}{\makecell{건물}} & \makecell{구 조} & \makecell{철골철근콘크리트조 (철근)콘크리트지붕} \\
\cline{2-3}
 & \makecell{용 도} & \makecell{업무시설} \\
\cline{2-3}
 & \makecell{건축면적} & \makecell{1,983.48m\textsuperscript{2}} \\
\cline{2-3}
 & \makecell{연면적} & \makecell{49,587m\textsuperscript{2}} \\
\cline{2-3}
 & \makecell{층 수} & \makecell{지상 20층/지하 5층} \\
\cline{2-3}
 & \makecell{사용승인} & \makecell{1999.12} \\
\cline{2-3}
 & \makecell{주 차} & \makecell{100대} \\
\cline{2-3}
 & \makecell{승강기} & \makecell{승객용 5(H사, 1,150kg, 90m/min, 15인)\newline{}비상용 1(H사, 750kg, 90m/min, 10인)} \\
\hline
\end{tabular}\par}

2. 대상물건 거래관련 자료 \par
{\small\centering
\begin{tabular}{|c|c|}
\hline
\makecell{거래예정금액} & \makecell{275,000,000,000원} \\
\hline
\makecell{거래조건} & \makecell{없 음} \\
\hline
\makecell{거래예정시점} & \makecell{2012.09.30} \\
\hline
\makecell{거래예정금액 지급조건} & \makecell{일시불(자기자본: 110,000,000,000원,\newline{}타인자본: 165,000,000,000원)} \\
\hline
\makecell{토지건물 배분예정금액} & \makecell{토지: 210,000,000,000원,\newline{}건물: 65,000,000,000원} \\
\hline
\makecell{오피스빌딩 하위시장} & \makecell{YS북부} \\
\hline
\end{tabular}\par}

\textbf{<자료 2> 시장임대료(Market Rent) 관련 자료} \par
1. 보증금, 연간임대료, 연간관리비, 고정경비 및 변동경비 상승률 \par
1) 1, 2년차: 5\% 또는 소비자물가지수(CPI; Consumer Price Index) 중 높은 율 \par
2) 3년차부터: 4\% 또는 CPI 중 높은 율 \par
2. 공실 및 대손충당금 \par
1) 가능총소득(PGI; Potential Gross Income)의 5\% 또는 PGI의 CPI \par
2) 일반경기의 회복지연으로 창업률이 낮아 공실률은 증가하고 있다. 따라서 대상물건의 감정평가시 공실 및 대손충당금 비율은 보수적인 측면을 고려하여 적용하여야 할 것으로 판단된다. \par
3. 보증금운용이율: 연 5\% \par
\textbf{<자료 3> 계약임대료(Contract Rent) 관련 자료} \par
1. 보증금, 연간임대료, 연간관리비, 고정경비 및 변동경비 상승률: 5\% 또는 CPI 중 낮은 율 \par
2. 공실 및 대손충당금 \par
1) PGI의 3.5\% \par
2) 대상물건은 양호한 임차인이 입주하고 있어 오피스빌딩 하위시장의 공실 및 대손충당금 비율에 비해 낮은 상태를 유지하고 있다. \par
3. 보증금운용이율: 연 5\% \par
\textbf{<자료 4> 공통자료} \par
1. 각종지표 \par
1) 2008년 금융위기 이후 일반경기의 본격적인 상승이 이루어지지 않고 있으며, CPI는 연간 3.5\% 상승이 예상된다. \par
2) 채권금리 등 \par
{\small\centering
\begin{tabular}{|c|c|c|c|}
\hline
\makecell{구 분} & \makecell{국고채(3년)} & \makecell{회사채(3년)} & \makecell{CD(91일)} \\
\hline
\makecell{\%} & \makecell{4.34} & \makecell{5.44} & \makecell{2.79} \\
\hline
\end{tabular}\par}

3) 벤치마크 투자수익률 \par
{\small\centering
\begin{tabular}{|c|c|c|c|c|}
\hline
\makecell{하위시장} & \makecell{SN북부} & \makecell{SN남부} & \makecell{YS북부} & \makecell{YS남부} \\
\hline
\makecell{투자수익률(\%)} & \makecell{6.0} & \makecell{6.2} & \makecell{7.0} & \makecell{6.5} \\
\hline
\end{tabular}\par}

2. 영업경비 \par
1) 고정경비: 연간관리비의 40\% \par
2) 변동경비: 연간관리비의 30\% \par
3) 대체충당금: 2년차에 100,000,000원, 4년차에 150,000,000원 설정 예정 \par
3. 할인율(Discount Rate) \par
1) 甲감정평가사는 할인율 결정방법을 (주)A감정평가법인의 감정평가심사위원회(이하 ‘위원회’)에 부의하여 결정하기로 하였다. \par
2) 동 위원회는 자본자산평가모델(CAPM; Capital Asset Pricing Model), 가중평균자본비용(WACC; Weighted Average Cost of Capital), 국고채금리에 일정률을 가산하여 구하는 방법 등을 종합 검토한 결과 본 감정평가에 적용할 할인율 결정방법은 WACC로 적용하는 것이 타당할 것 같다고 甲감정평가사에게 권고하였고, 甲감정평가사는 이를 수용하였다. \par
3) 甲감정평가사가 오피스빌딩 하위시장에서 조사한 자기자본수익률은 6.50\%, 타인자본수익률은 5.67\%이다. \par
4) 甲감정평가사가 대상물건과 관련하여 조사한 자기자본수익률은 6.25\%, 타인자본수익률은 5.00\%이다. \par
5) IRR, WACC는 소수점 넷째자리 이하 절사한다. \par
4. 기출환원이율(Terminal Cap Rate) \par
1) 甲감정평가사는 기출환원이율도 위원회에 부의하여 결정하기로 하였다. \par
2) 위원회는 향후 오피스빌딩시장의 가격변동률이 하락될 것으로 예상되어 기출환원이율의 결정은 보수적인 입장을 취하는 것이 합리적이라는 판단을 하였다. \par
3) 그래서 위원회는 甲감정평가사에게 기출환원이율은 결정된 할인율에 0.5\%p를 가산하여 적용하는 것이 타당할 것이라는 권고를 하였다. 甲감정평가사도 오피스빌딩 시장이 부동산경기 및 일반경기 침체로 인해 하락할 것으로 판단하여 위원회의 권고안을 받아들이기로 하였다. \par
5. 보유기간 및 복귀가격 \par
1) 보유기간은 5년으로 하고, 재매도비용은 공인중개사수수료 및 기타비용 등을 고려할 때 2\%로 한다. \par
2) 복귀가격 결정을 위한 PGI 등은 5년차의 PGI 등에 연간 임대료, 관리비 등의 상승률을 적용하여 6년차의 순영업소득(NOI; Net Operating Income)을 기준으로 결정한다. \par
6. 할인율과 기별 계수 \par
{\small\centering
\begin{tabular}{|c|c|c|c|c|c|}
\hline
\makecell{구분} & \makecell{1기} & \makecell{2기} & \makecell{3기} & \makecell{4기} & \makecell{5기} \\
\hline
\makecell{2.79\%} & \makecell{0.973} & \makecell{0.946} & \makecell{0.921} & \makecell{0.896} & \makecell{0.871} \\
\hline
\makecell{4.34\%} & \makecell{0.958} & \makecell{0.919} & \makecell{0.880} & \makecell{0.844} & \makecell{0.809} \\
\hline
\makecell{5.00\%} & \makecell{0.952} & \makecell{0.907} & \makecell{0.864} & \makecell{0.823} & \makecell{0.784} \\
\hline
\makecell{5.44\%} & \makecell{0.948} & \makecell{0.899} & \makecell{0.853} & \makecell{0.809} & \makecell{0.767} \\
\hline
\makecell{5.50\%} & \makecell{0.948} & \makecell{0.898} & \makecell{0.852} & \makecell{0.807} & \makecell{0.765} \\
\hline
\makecell{5.80\%} & \makecell{0.945} & \makecell{0.893} & \makecell{0.844} & \makecell{0.798} & \makecell{0.754} \\
\hline
\makecell{5.94\%} & \makecell{0.944} & \makecell{0.891} & \makecell{0.841} & \makecell{0.794} & \makecell{0.749} \\
\hline
\makecell{6.00\%} & \makecell{0.943} & \makecell{0.890} & \makecell{0.840} & \makecell{0.792} & \makecell{0.747} \\
\hline
\makecell{6.50\%} & \makecell{0.939} & \makecell{0.882} & \makecell{0.828} & \makecell{0.777} & \makecell{0.730} \\
\hline
\makecell{7.00\%} & \makecell{0.935} & \makecell{0.873} & \makecell{0.816} & \makecell{0.763} & \makecell{0.713} \\
\hline
\makecell{7.50\%} & \makecell{0.930} & \makecell{0.865} & \makecell{0.805} & \makecell{0.749} & \makecell{0.697} \\
\hline
\makecell{8.00\%} & \makecell{0.926} & \makecell{0.857} & \makecell{0.794} & \makecell{0.735} & \makecell{0.681} \\
\hline
\makecell{8.50\%} & \makecell{0.922} & \makecell{0.849} & \makecell{0.783} & \makecell{0.722} & \makecell{0.665} \\
\hline
\makecell{9.00\%} & \makecell{0.917} & \makecell{0.842} & \makecell{0.772} & \makecell{0.708} & \makecell{0.650} \\
\hline
\makecell{9.50\%} & \makecell{0.913} & \makecell{0.834} & \makecell{0.762} & \makecell{0.696} & \makecell{0.635} \\
\hline
\makecell{10.00\%} & \makecell{0.909} & \makecell{0.826} & \makecell{0.751} & \makecell{0.683} & \makecell{0.621} \\
\hline
\end{tabular}\par}

7. 임대사례 \par
{\small\centering
\begin{tabular}{|c|c|c|c|c|}
\hline
\makecell{구분} & \makecell{임대사례 1} & \makecell{임대사례 2} & \makecell{임대사례 3} & \makecell{대상물건} \\
\hline
\makecell{소재지} & \makecell{서울특별시\newline{}G구 Y동} & \makecell{서울특별시\newline{}G구 Y동} & \makecell{서울특별시\newline{}G구 Y동} & \makecell{서울특별시\newline{}G구 Y동} \\
\hline
\makecell{건물명} & \makecell{××빌딩} & \makecell{gg빌딩} & \makecell{zz빌딩} & \makecell{OO빌딩} \\
\hline
\makecell{층(지상/지하)} & \makecell{30F/B3} & \makecell{20F/B6} & \makecell{25F/B6} & \makecell{20F/B5} \\
\hline
\makecell{구조} & \makecell{철골철근\newline{}콘크리트} & \makecell{철골철근\newline{}콘크리트} & \makecell{철골철근\newline{}콘크리트} & \makecell{철골철근\newline{}콘크리트} \\
\hline
\makecell{건물연면적(m\textsuperscript{2})} & \makecell{66,000} & \makecell{46,280} & \makecell{59,985} & \makecell{49,587} \\
\hline
\makecell{토지면적(m\textsuperscript{2})} & \makecell{3,333} & \makecell{2,966} & \makecell{3,225} & \makecell{2,833} \\
\hline
\makecell{전용률(\%)} & \makecell{53.1} & \makecell{48.4} & \makecell{52.1} & \makecell{49.8} \\
\hline
\makecell{사용승인(년)} & \makecell{1991} & \makecell{2000} & \makecell{2010} & \makecell{1999} \\
\hline
\makecell{오피스빌딩 하위시장} & \makecell{YS북부} & \makecell{YS북부} & \makecell{YS남부} & \makecell{YS북부} \\
\hline
\makecell{보증금(원/m\textsuperscript{2})} & \makecell{180,000} & \makecell{210,000} & \makecell{300,000} & \makecell{240,000} \\
\hline
\makecell{월임대료(원/m\textsuperscript{2})} & \makecell{18,000} & \makecell{21,000} & \makecell{30,000} & \makecell{24,000} \\
\hline
\makecell{월관리비(원/m\textsuperscript{2})} & \makecell{8,000} & \makecell{8,000} & \makecell{11,000} & \makecell{10,000} \\
\hline
\end{tabular}\par}

\end{longquestion}
